% =====================================================================
%  Bloc signature / marque — HAMADY CISSE
%  Réutilisable : copier ce fichier dans un projet et faire
%      % =====================================================================
%  Bloc signature / marque — HAMADY CISSE
%  Réutilisable : copier ce fichier dans un projet et faire
%      % =====================================================================
%  Bloc signature / marque — HAMADY CISSE
%  Réutilisable : copier ce fichier dans un projet et faire
%      % =====================================================================
%  Bloc signature / marque — HAMADY CISSE
%  Réutilisable : copier ce fichier dans un projet et faire
%      \input{signature.tex}      (dans le préambule)
%  puis appeler  \signatureHamady  là où la signature doit apparaître.
%  Dépendances : xcolor (déjà chargé dans les rapports du mentorat).
% =====================================================================

\definecolor{sigblue}{HTML}{1E70BF}

% Empreinte GPG (fingerprint) — à remplacer après génération de la clé.
\newcommand{\gpgEmpreinte}{A6F4 0DFD EE3E B7AE 80B3\ 6E77 361C 9E02 B790 303E}

\newcommand{\signatureHamady}{%
  \par\vspace{1.2em}%
  \noindent{\color{sigblue}\rule{\linewidth}{0.6pt}}\par\vspace{0.4em}%
  \noindent{\bfseries\large HAMADY CISSE}\par\vspace{0.15em}%
  \noindent{\small Diplômé en Cryptographie et Sécurité de l'Information}\par\vspace{0.35em}%
  \noindent{\footnotesize\ttfamily Courriel : hamadycisse7540@gmail.com \quad GitHub : hamadycisse7540-hash}\par\vspace{0.1em}%
  \noindent{\footnotesize\ttfamily Empreinte GPG : \gpgEmpreinte}\par\vspace{0.3em}%
  \noindent{\scriptsize\itshape Document signé numériquement (GPG). Vérifiez son authenticité et son intégrité avec la clé publique ci-dessus.}\par\vspace{0.4em}%
  \noindent{\color{sigblue}\rule{\linewidth}{0.6pt}}\par
}
      (dans le préambule)
%  puis appeler  \signatureHamady  là où la signature doit apparaître.
%  Dépendances : xcolor (déjà chargé dans les rapports du mentorat).
% =====================================================================

\definecolor{sigblue}{HTML}{1E70BF}

% Empreinte GPG (fingerprint) — à remplacer après génération de la clé.
\newcommand{\gpgEmpreinte}{A6F4 0DFD EE3E B7AE 80B3\ 6E77 361C 9E02 B790 303E}

\newcommand{\signatureHamady}{%
  \par\vspace{1.2em}%
  \noindent{\color{sigblue}\rule{\linewidth}{0.6pt}}\par\vspace{0.4em}%
  \noindent{\bfseries\large HAMADY CISSE}\par\vspace{0.15em}%
  \noindent{\small Diplômé en Cryptographie et Sécurité de l'Information}\par\vspace{0.35em}%
  \noindent{\footnotesize\ttfamily Courriel : hamadycisse7540@gmail.com \quad GitHub : hamadycisse7540-hash}\par\vspace{0.1em}%
  \noindent{\footnotesize\ttfamily Empreinte GPG : \gpgEmpreinte}\par\vspace{0.3em}%
  \noindent{\scriptsize\itshape Document signé numériquement (GPG). Vérifiez son authenticité et son intégrité avec la clé publique ci-dessus.}\par\vspace{0.4em}%
  \noindent{\color{sigblue}\rule{\linewidth}{0.6pt}}\par
}
      (dans le préambule)
%  puis appeler  \signatureHamady  là où la signature doit apparaître.
%  Dépendances : xcolor (déjà chargé dans les rapports du mentorat).
% =====================================================================

\definecolor{sigblue}{HTML}{1E70BF}

% Empreinte GPG (fingerprint) — à remplacer après génération de la clé.
\newcommand{\gpgEmpreinte}{A6F4 0DFD EE3E B7AE 80B3\ 6E77 361C 9E02 B790 303E}

\newcommand{\signatureHamady}{%
  \par\vspace{1.2em}%
  \noindent{\color{sigblue}\rule{\linewidth}{0.6pt}}\par\vspace{0.4em}%
  \noindent{\bfseries\large HAMADY CISSE}\par\vspace{0.15em}%
  \noindent{\small Diplômé en Cryptographie et Sécurité de l'Information}\par\vspace{0.35em}%
  \noindent{\footnotesize\ttfamily Courriel : hamadycisse7540@gmail.com \quad GitHub : hamadycisse7540-hash}\par\vspace{0.1em}%
  \noindent{\footnotesize\ttfamily Empreinte GPG : \gpgEmpreinte}\par\vspace{0.3em}%
  \noindent{\scriptsize\itshape Document signé numériquement (GPG). Vérifiez son authenticité et son intégrité avec la clé publique ci-dessus.}\par\vspace{0.4em}%
  \noindent{\color{sigblue}\rule{\linewidth}{0.6pt}}\par
}
      (dans le préambule)
%  puis appeler  \signatureHamady  là où la signature doit apparaître.
%  Dépendances : xcolor (déjà chargé dans les rapports du mentorat).
% =====================================================================

\definecolor{sigblue}{HTML}{1E70BF}

% Empreinte GPG (fingerprint) — à remplacer après génération de la clé.
\newcommand{\gpgEmpreinte}{A6F4 0DFD EE3E B7AE 80B3\ 6E77 361C 9E02 B790 303E}

\newcommand{\signatureHamady}{%
  \par\vspace{1.2em}%
  \noindent{\color{sigblue}\rule{\linewidth}{0.6pt}}\par\vspace{0.4em}%
  \noindent{\bfseries\large HAMADY CISSE}\par\vspace{0.15em}%
  \noindent{\small Diplômé en Cryptographie et Sécurité de l'Information}\par\vspace{0.35em}%
  \noindent{\footnotesize\ttfamily Courriel : hamadycisse7540@gmail.com \quad GitHub : hamadycisse7540-hash}\par\vspace{0.1em}%
  \noindent{\footnotesize\ttfamily Empreinte GPG : \gpgEmpreinte}\par\vspace{0.3em}%
  \noindent{\scriptsize\itshape Document signé numériquement (GPG). Vérifiez son authenticité et son intégrité avec la clé publique ci-dessus.}\par\vspace{0.4em}%
  \noindent{\color{sigblue}\rule{\linewidth}{0.6pt}}\par
}
